\chapter*{使用说明}
\addcontentsline{toc}{chapter}{使用说明}
\markboth{使用说明}{使用说明}
\begin{enumerate}
\item \key{定位}：本册只收考场上要查的内容——概念定义、模型与损失函数、优化与更新公式、算法步骤、评价指标、方法对比、Python 代码模板与常用函数、常见报错与易错点、英文一句话表述。不含习题。
\item \key{结构}：第一部分（第 1–3 章）机器学习基础：基本概念与三要素、线性模型（Logistic/Softmax/感知器/SVM）、前馈神经网络与反向传播；第二部分（第 4–5 章）模型评估、经典监督学习算法、数据预处理、聚类、降维、关联规则与经济金融应用；第三部分（第 6–8 章）Python 爬虫（requests、正则、Selenium、巨潮资讯、PDF 解析）与文本分析（分词、TF-IDF、情感分析、文本分类、LDA）；附录为 Python 速查、常见报错与英文答题句式、中英术语速查。
\item \key{怎么查}：目录细到小知识点；页眉显示章名（左页）与节名（右页）。英文术语查附录 C；代码函数查附录 A；报错与概念辨析查附录 B。
\item \key{标记}：公式推导右侧灰色〔……〕为该步理由；\tagext\ 为拓展内容（超出核心范围但可能考到）；“易错”条目列出最常见的错误。
\item \key{代码约定}：代码按当前库版本书写（Python 3、requests、re、Selenium 4、pdfplumber、jieba、scikit-learn、gensim）。
\end{enumerate}
